\chapter[TINJAUAN PUSTAKA]{\\ TINJAUAN PUSTAKA}
\section{Landasan Teoritis}
[ISI KERANGKA TEORI DAN DEFINISI KONSTRUK.]

\section{Tinjauan Hasil Penelitian Terkait}
[ISI SINTESIS KRITIS STUDI TERDAHULU DAN SITASI TERVERIFIKASI.]

\subsection{Kondisi Mutakhir Penelitian (\textit{State of the Art})}
[ISI PETA TEMUAN DAN PENDEKATAN RISET.]

\subsection{Kesenjangan Penelitian}
[ISI GAP EKSPLISIT DAN GAP HASIL ANALISIS.]

\subsection{Posisi Penelitian}
[ISI KONTRIBUSI YANG DIUSULKAN DALAM KONTEKS LITERATUR.]
